% ============================================================
%  第4章  软件设计与实现
% ============================================================
\chapter{软件设计与实现}

本章详细阐述头戴式眼部康复装置的软件系统设计，包括整体架构、姿态解算算法、五种临床训练模式的实现逻辑、语音交互机制以及安全保护策略。系统以STM32H723VGT6为主控核心，基于HAL库和CubeMX框架进行开发，软件主体采用前后台架构，以10ms周期主循环配合中断驱动外设，实现实时姿态跟踪与多模式康复训练。

\section{软件整体架构}

系统的软件架构遵循“感知—决策—执行—反馈”的闭环设计思想，软件整体架构框图如图\ref{fig:software_architecture}所示。感知层通过SPI总线采集BMI088惯性传感器数据，经姿态解算获得头部的俯仰、偏航与侧倾欧拉角；同时通过UART接收CI1302语音模块识别到的命令词。决策层的核心是一个九状态系统状态机，根据当前状态和外部事件（语音命令、姿态报警、按键输入）进行状态转移和训练逻辑控制。执行层驱动舵机云台完成激光点的空间定位，并控制激光模块、LED指示灯和蜂鸣器输出视觉和听觉刺激。反馈层则通过语音播报进行结果评价与动作指导，通过WS2812状态灯提供持续的生物反馈。

\begin{figure}[h!]
\centering
\includegraphics[width=1.0\textwidth]{sw_arch.drawio.pdf}
\caption{系统软件架构框图}
\label{fig:software_architecture}
\end{figure}

\subsection{主循环结构}

系统的主循环采用10ms固定节拍作为调度基准，其执行流程参见图\ref{fig:main_loop}。每轮循环按以下顺序逐层执行：率先通过欧拉角读取函数从惯性传感器采集最新原始数据并刷新欧拉角估计；将更新后的姿态信息传至头部姿态分析模块以重算头稳等指标；随后轮询扫描PA2按键的触发状态；最后依当前状态机的取值将控制权分发至对应的训练状态处理例程，并在循环末尾插入一段10ms固定延时以保证迭代周期的严格等间隔。\begin{figure}[h!]
\centering
\includegraphics[width=1.0\textwidth]{figures/main_loop.drawio.pdf}
\caption{主循环流程图}
\label{fig:main_loop}
\end{figure}

为了提升系统的可观测性和调试效率，主循环中嵌入了一个基于UART的VFo（Volatile Flow）协议输出模块。该模块以100ms为周期，通过USART1向PC端上位机发送结构化数据帧，包含力加速度、角速度、欧拉角、头稳指标、系统状态、最后语音命令等关键运行时参数。这种在线监测方式为开发阶段的问题定位和debug提供了重要支持。

\subsection{系统状态机设计}

系统全部运行逻辑由一套九状态状态机集中管理，状态间的转移关系已在2.2.3节的图\ref{fig:train_flow}中呈现。与早期原型中的五状态方案对比，新增的四个状态分别承担了双训练模式选择、模式进入语音预告、每次训练前的引导播报以及完整训练链路中各模式间的确认衔接功能，下文给出各状态的详细定义。

(1) 系统模式选择状态：系统初始化完成后所处的第一个调度状态。在该状态下，所有刺激输出通道——激光、聚焦LED和舵机——均保持关闭，语音模块播报引导语"按一下按键进入完整训练，按两下进入自选模式"。患者借助PA2按键完成模式分配：若在规定双击窗口（600ms）内未捕获到第二次按键下降沿，则判定为单击操作，对应进入完整训练流程（APP\_FLOW\_FULL），系统按A→B→C→D→E的预设顺序依次串联全部五种模式；若两次按下间隔短于600ms则识别为双击，对应进入自选训练流程（APP\_FLOW\_CUSTOM），此后由患者通过语音命令逐一挑选所欲执行的训练项目。日常运行中如需在两种模态间切换，患者可在任意时间说出"模式切换"命令词直接重回此状态。

(2) 系统模式提示状态：模式选定后即刻进入的过渡态。系统根据所选是完整模式还是自选模式播报不同的确认提示语，静默2.5s后自动跳转——完整模式从模式A的入口开始训练链路，自选模式则进入空闲状态挂起等待语音指令。

(3) 系统空闲等待状态：自选训练流程下的核心驻留状态。系统在此状态下不执行任何训练动作，仅保持WS2812指示灯以绿色呼吸灯节律闪烁以表明就绪，同时对语音模块唤醒后识别的有效命令词（如"扫视训练""注视稳定训练""标定模式"等）保持全时监听，收到匹配命令后立即转至姿态校准流程。

(4) 系统姿态校准状态：接到有效训练命令后第一步进入的前置准备态。语音播报"请保持头部静止"，延迟2.5s后对陀螺仪三轴采集800个连续静止采样点并求取均值，以此作为当前运行环境下的零偏基准写入全局偏置寄存器。校准一经完成即自动转至系统训练提示状态。

(5) 系统训练提示状态：紧邻每次具体训练开始之前的预告阶段。系统结合当前训练模式的类别播报对应的专用引导语音——例如注视训练播报"请注视光点保持头部稳定"、扫视训练播报"请快速找到光点并按按键"、追踪训练播报"请用眼睛跟随光点移动"等。为完整覆盖较慢语速播报的全部内容，系统在提示播放后等待4s再转入系统训练执行状态正式启动训练。

(6) 系统训练执行状态：训练过程的主体运行阶段。每帧循环内先执行姿态安全检查，随后将控制权交予当前模式对应的训练函数执行具体方案。五种训练模式在该状态下共享统一的暂停/恢复接口与语音反馈通道，确保模式间调度逻辑的一致性。

(7) 系统反馈评价状态：任一训练项目执行完毕后自动切至此态。系统汇集训练过程中的正确率、反应时间、头稳指标等量化记录，依据预设的分级标准给出综合评价并通过语音播报。在完整训练流程下，反馈结束后系统进入下一模式确认状态等候患者按键衔接；自选流程下则返回系统空闲等待状态继续监听新的语音选训命令。

(8) 系统下一模式确认状态：仅在完整训练流程中被激活。当一项训练完结且反馈评价播报完毕后即进入此态，语音提示"按键确认进入下一个模式"。患者按下PA2后模式指针前移一位（例：注视→扫视），流程重置至系统姿态校准状态开始下一轮的校准—提示—执行—反馈链。系统在确认前保持静默，不设超时自动跳转，完全由患者自主把控节奏。

(9) 系统安全暂停状态：训练执行中若检测到任一轴姿态超限或收到语音"暂停训练"命令时触发。激光与聚焦LED立即熄灭，舵机执行回中动作，蜂鸣器发出短促告警脉冲，同时语音模块播放对应警告提示语。系统对两种暂停触发源做了区分处理：由姿态超限引起的硬件保护暂停在姿态连续正常满3s后自动恢复且不丢失训练进度；由语音命令发起的监听暂停则在5s监听窗口内持续扫描后续控制指令（"继续训练""重新开始"等），超时无后续指令同样自动恢复当前训练。

为方便运行状态的离线回溯，系统在每次状态发生转移时均记录转移事件类型及关联的时间戳，并通过前述VFo协议输出至PC端，为调试和性能剖析提供完整的运行时日志轨迹。
\section{BMI088驱动与姿态解算}

\subsection{SPI通信初始化}

BMI088属于博世Sensortec公司推出的工业级6轴惯性传感器，封装内集成了相互电气隔离的16位三轴加速度计和16位三轴陀螺仪~\cite{bmi088_ds}。两颗子传感芯体通过各自独立的片选控制线（加速度计由PC0即ACC\_CS选通，陀螺仪由PC3即GYRO\_CS选通）挂接在同一条SPI总线下。主控端启用STM32H723的SPI2外设以全双工方式~\cite{stm32h7_rm}与BMI088交换数据，两枚传感芯片的中断就绪信号则分别经PE10和PE12引脚接入MCU的中断输入通道。

上电初始化序列按以下步骤逐项推进：先分别向加速度计和陀螺仪发送硬件复位命令；随后回读两者的芯片ID寄存器比对期望值以确认物理通信链路已建立；再对加速度计一侧配置量程为$\pm$3g、输出数据率为800Hz、工作模式为Normal并令INT1引脚以开漏模式输出，对陀螺仪一侧配置量程为$\pm$2000dps、带宽取1000Hz/116Hz组合、同样设为Normal模式且由INT3引脚输出。所有寄存器的写入结果在配置流程末尾均经过回读比对校验。

\subsection{坐标系定义与数据流向}

BMI088芯片在DM-02核心板上的实际焊接方向为：芯片X敏感轴沿板面指向USB接口一侧（竖直朝上），Y敏感轴指向XT30电源输入接口一侧（水平向左），Z敏感轴的方向根据右手定则判定为垂直于芯片外表面指向前方（即背离佩戴者头部）。在这一安装姿态下，BMI088自身的三轴敏感方向与佩戴者头部的自然转动轴之间的映射关系统一列于表\ref{tab:axis_mapping}，其在开发板上的欧拉角物理方向示意如图\ref{fig:euler_angles}所示。

\begin{table}[h!]
\centering
\caption{BMI088传感器轴与佩戴姿态的旋转对应关系}
\label{tab:axis_mapping}
\begin{tabular}{p{2.5cm} p{3.5cm} p{5.5cm}}
\toprule[1.5pt]
\textbf{BMI088 轴} & \textbf{物理方向（USB口朝上佩戴）} & \textbf{旋转对应} \\
\midrule[1pt]
X轴 & 垂直向上 & 绕垂直轴的旋转（左右转头） \\
Y轴 & 水平向左 & 绕左右轴的旋转（点头仰头） \\
Z轴 & 垂直芯片表面指向前方 & 绕前后轴的旋转（侧倾/耳朵靠肩） \\
\bottomrule[1.5pt]
\end{tabular}
\end{table}

\begin{figure}[htbp]
\centering
\includegraphics[width=0.98\textwidth]{euler.drawio.pdf}
\caption{欧拉角物理方向示意图}
\label{fig:euler_angles}
\end{figure}

上述映射关系恰好与工程实践中惯用的欧拉角物理定义完全吻合——横滚角（roll）的轴对应左右转头的旋转、俯仰角（pitch）的轴对应点头/仰头的旋转、偏航角（yaw）的轴对应侧倾运动——因此系统无需在传感器输出与姿态算法输入之间插入额外的坐标变换矩阵，直接将BMI088三轴原始测量值送入姿态解算链路。这种直通式数据流既精简了代码逻辑，也从根本上规避了因手写重映射矩阵中的轴序符号错误而诱发的长期调试问题。

\subsection{姿态估计与数据融合}

系统当前的姿态解算管线已全面迁移至基于四元数的融合框架。与早期开发阶段为每个轴单独设计姿态策略（卡尔曼滤波用于roll轴、纯陀螺仪积分驱动yaw轴、加速度计正切反算求解pitch轴）的方案不同，当前版本采用四元数作为统一的姿态表达载体，在一个共同的数学空间内完整处理三个旋转自由度的估计问题。

四元数可视为复数概念向四维空间的自然推广~\cite{madgwick2011orientation,han2011novel}，通常写作$\mathbf{q} = q_0 + q_1\mathbf{i} + q_2\mathbf{j} + q_3\mathbf{k}$，其中标量分量$q_0$为实部，矢量分量$(q_1, q_2, q_3)$构成虚部，单位基$\mathbf{i},\mathbf{j},\mathbf{k}$服从关系$\mathbf{i}^2=\mathbf{j}^2=\mathbf{k}^2=\mathbf{ijk}=-1$。对于三维旋转的描述而言，单位四元数（即满足$|\mathbf{q}|=1$）能够将绕单位轴$\mathbf{u}$旋转角度$\theta$的任意空间转动紧凑地表示为$\mathbf{q} = [\cos(\theta/2),\; \mathbf{u}\sin(\theta/2)]$。

联系陀螺仪测量与姿态变化的动力学方程构成了四元数姿态更新的数学核心。记BMI088陀螺仪输出的三轴角速度向量为$\boldsymbol{\omega}_s = [\omega_{sx}, \omega_{sy}, \omega_{sz}]$，将其升维为纯四元数形式$\boldsymbol{\omega} = (0, \omega_{sx}, \omega_{sy}, \omega_{sz})$后，姿态四元数的时间导数由下式\eqref{eq:quat_diff}给出：

\begin{equation}
\dot{\mathbf{q}} = \frac{1}{2} \boldsymbol{\omega} \otimes \mathbf{q}
\label{eq:quat_diff}
\end{equation}

上式中$\otimes$为四元数乘法算子。该微分方程的工程解释可以概括为：每一帧内，陀螺仪测出的绕三轴旋转的瞬时角速度（rad/s）乘以采样间隔$\Delta t$得出该帧内的微元旋转角，将其封装为增量四元数$\Delta\mathbf{q}$后左乘至当前时刻的姿态四元数，即完成姿态的一步前向传播，如式\eqref{eq:quat_integrate}所示：

\begin{equation}
\mathbf{q}_{t+1} = \mathbf{q}_t + \dot{\mathbf{q}}_t \cdot \Delta t
\label{eq:quat_integrate}
\end{equation}

这一单次四元数乘法的数学性质已经将三轴旋转之间的顺序耦合关系完整地纳入了计算过程，由此彻底避免了欧拉角方案中因逐轴积分且旋转次序不可交换所造成的坐标转换累积漂移。

本项目选择四元数而舍弃传统欧拉角的决策主要基于以下三点考量。

（1）规避万向锁奇异。欧拉角表示在$pitch = \pm 90^\circ$附近遭遇万向锁——此时横滚和偏航两轴退化至同一旋转平面，系统实际丧失了一个独立自由度。患者在佩戴装置后的日常头部活动中完全可能不经意接近该奇点（例如大幅弯腰低头寻找掉落物品的情形），一旦进入此区域，欧拉角输出将产生不连续的跳变。四元数作为全局无奇点的姿态参数化方法，在任意方位角下均维持姿态演化的光滑性与连续性。

（2）角速度的一体化积分。BMI088陀螺仪以固定频率向外推送三轴角速度采样。若沿用欧拉角方案，则三轴角速度须依次按Roll→Pitch→Yaw的既定序列分步积分，而旋转顺序本身对最终姿态的精度就有不可忽略的影响。相比之下，借助式\eqref{eq:quat_diff}这一统一的微分方程，只需一次四元数乘法和一次归一化即可推进姿态预测的全部计算，运算结构紧凑且物理背景清晰。

（3）加速度计修正在四元数域内天然稳定。将由加速度计分量推算出的观测roll和pitch（式\eqref{eq:roll_acc}、式\eqref{eq:pitch_acc}）包装为重力参考四元数$\mathbf{q}_{\text{acc}}$，再与陀螺预测量$\mathbf{q}_{\text{pred}}$执行互补融合：

\begin{equation}
\mathbf{q}_{\text{corrected}} = \mathbf{q}_{\text{pred}} + K \cdot (\mathbf{q}_{\text{acc}} - \mathbf{q}_{\text{pred}})
\label{eq:quat_correct}
\end{equation}

修正增益$K$取默认值2.8。融合结果经归一化后直接折算为三轴欧拉角输出。由于四元数空间内的加法与标量乘法天然平滑，修正过程中不会出现欧拉角域内常见的$180^\circ$与$-180^\circ$之间的边界翻转问题，保证了姿态估计序列的时间一致性。

在具体的运行时序上，姿态预测环节读取陀螺仪的三轴角速度构造增量四元数，通过式\eqref{eq:quat_diff}和式\eqref{eq:quat_integrate}将当前姿态向前推进一帧。姿态校正环节以加速度计归一化后的重力方向向量作为绝对空间参照，借助式\eqref{eq:roll_acc}和式\eqref{eq:pitch_acc}反算出观测俯仰和横滚角度，再经式\eqref{eq:quat_correct}完成融合。修正后的四元数归一化之后即转换为三轴欧拉角，作为头部姿态分析与安全监测等全部下游模块的唯一姿态信息来源。

\begin{equation}
roll = \text{atan2}(a_{sy},\; a_{sz}) \cdot \frac{180}{\pi}
\label{eq:roll_acc}
\end{equation}

\begin{equation}
pitch = \text{atan2}(-a_{sx},\; \sqrt{a_{sy}^2 + a_{sz}^2}) \cdot \frac{180}{\pi}
\label{eq:pitch_acc}
\end{equation}

在三种姿态角分量中，横滚角因重力在加速度计Y-Z平面上的正交投影而受益于直接物理约束，在静态与准静态条件下的估计精度很高。俯仰角由加速度计X轴分量与Y-Z平面合成矢量联合反算，当头部未处于极端俯仰姿态时（$|pitch| < 60^\circ$）同样具有良好的可信度。偏航角因缺乏磁力计提供的地磁绝对方位基准，仅能依靠陀螺仪Z轴角速度的逐帧积分来维持，长时间运行下不可避免会积累漂移。通过在初始化阶段采集800个静止样本对三轴陀螺仪零偏进行校准，单次典型训练周期（$15\sim50$s）内yaw累积漂移量可控制在$2^\circ$以内，足以满足康复训练应用对侧倾感知的精度要求。

系统在静止状态下对陀螺仪各轴的偏置执行在线自适应跟踪：当陀螺仪输出合角速度降至0.08\,dps以下时，以指数平滑系数$\alpha=0.0025$将当前轴的偏置估计朝最新陀螺仪读数方向缓慢逼近，从而持续补偿环境温度变化等因素引起的缓变漂移。通信异常防护层面，设置了3\,ms的超时回退机制——在偶然出现SPI总线竞争或DMA缓存未就绪的工况下，姿态输出自动退避为前帧有效值的延续，确保三轴欧拉角序列中不出现NaN或阶跃跳点。
\section{头部姿态分析与安全监测}

\subsection{头稳指标计算}

头部的稳态维持能力是追踪康复成效的核心量化依据，在注视稳定性训练中它以最直接的方式表征了患者抑制头部额外运动的能力。系统构建了一个滑动窗口方差估计器来量化该指标：维护容量为100帧（覆盖1s的历史时间跨度）的环形缓冲队列，其中分别容纳俯仰角与横滚角的连续100个离散采样。每帧迭代中同步刷新两通道样本集的方差，取方差之和的平方根作为当前瞬时头稳值：

\begin{equation}
\text{stability} = \sqrt{\text{Var}(\theta) + \text{Var}(\phi)}
\label{eq:stability}
\end{equation}

此指标的量纲为度（°），数值越小指示头部越趋近于静止。当该指标连续10帧以上维持在3°的门限上方时，系统即判定佩戴者存在病理性头抖倾向并激活安全暂停机制。

\subsection{代偿性转头检测}

脑卒中后罹患偏盲的患者在认知代偿机制驱动下，常常不自觉地以躯干或颈部的整体旋转来替代眼球的快速扫视运动。从短期看这种代偿策略似乎提高了目标搜索效率，但从中长期神经可塑性角度衡量，它恰恰抑制了受损视路和眼外肌群的主动重激活过程。鉴于此，本系统将过度的代偿性转头行为明确定义为训练过程中需要识别并纠正的不良模式。

依据头部坐标系约定，横滚角构成反映左右转头幅度的直接量度。系统设定代偿判据阈值为20°：一旦横滚角的瞬时绝对值跨越该界值，标志位$is\_compensatory$随即置位，当前训练模式据此触发双通道警示——语音播报"请用眼睛跟踪，不要转头"，WS2812指示灯同步从绿色呼吸切换为红色高速闪烁。

\subsection{安全倾斜检测}

除代偿转头外，训练持续期间还应防范因患者疲劳或坐姿不当而引发的躯干过度前倾/后仰。此类姿态不仅存在跌倒致伤的现实风险，同时也会使激光光斑在墙面上的投影偏离预期的训练作业区域，削弱训练效果。系统对俯仰角设置的安全门限为30°：当绝对值超越该值即标记$safety\_alert$，系统即刻进入暂停态——激光立即熄灭、双轴舵机执行回中复位、蜂鸣器响短促告警、语音模块同步播报"训练暂停，请调整坐姿"。

\subsection{姿势识别}

在训练实施全过程中，系统对佩戴者当前的头部姿势执行持续分类，用以辅助评估训练姿态的规范程度。分类判定规则基于当前帧俯仰角与横滚角的相互大小逻辑：当俯仰角绝对值超过20°且在幅值上压倒横滚角时，正俯仰对应仰头类型（POSTURE\_UP），负俯仰对应低头类型（POSTURE\_DOWN）；反之当横滚角绝对值超过15°且在数值上强于俯仰角时，正横滚映射为左偏姿态（POSTURE\_LEFT），负横滚映射为右偏姿态（POSTURE\_RIGHT）；两角分量均未超出各自门限的情形则视为正中位（POSTURE\_CENTER）。
\section{五种训练模式的设计与实现}

\subsection{训练状态机总体设计}

五种训练模式统一由系统状态机调度管理。系统支持两种工作模式——完整训练模式（Full Mode）和自选训练模式（Custom Mode），由患者在系统模式选择状态下通过PA2按键选择。

完整训练模式的运行流程如图\ref{fig:full_mode_flow}所示：

\begin{figure}[h!]
\centering
\includegraphics[width=1.0\textwidth]{Full_Mode.drawio.pdf}
\caption{完整训练模式}
\label{fig:full_mode_flow}
\end{figure}

五种模式按A→B→C→D→E的固定顺序依次执行，每完成一个模式后需患者按键确认方可进入下一模式，五次完成后自动回到系统模式选择状态。

自选训练模式的运行流程如图\ref{fig:custom_mode_flow}所示:

\begin{figure}[h!]
\centering
\includegraphics[width=1.0\textwidth]{Custom_Mode.drawio.pdf}
\caption{自选训练模式}
\label{fig:custom_mode_flow}
\end{figure}

模式命令词被设计为可在IDLE\_VOICE、TRAIN、PAUSE和FEEDBACK四个状态下跨状态响应，即训练执行中或暂停状态下均可直接切换训练模式而无需先回到空闲状态。

每项训练正式开始前，系统均会先进入系统训练提示状态播放该模式专用的开始提示语音（如“注视训练开始，请注视光点保持头部稳定”），给予老年患者明确的听觉预告和准备时间，4s后自动进入训练。

所有训练模式共享一套统一的患者反馈接口：舵机控制通过调用舵机角度设置函数（axis, angle）实现，该函数同时记录当前舵机目标角度，供暂停恢复时使用；激光和LED的控制同样通过应用层封装函数App调用激光控制函数和App调用聚焦LED控制函数进行，确保状态一致性。

\subsection{模式A：注视稳定性训练}

注视稳定性训练的核心任务是对患者维持头部静止不动的基本控制力进行测试与循序提升。在训练启动阶段，舵机云台首先将激光投射至视野正前方由标定确定的预设坐标点，随后患者需在15s倒计时区间内持续以视线锁住该固定光点，同时有意识地约束头部的多余摇晃。

系统以帧为单位滚动计算头稳指标：一旦该值超越5°的门限，语音通道即刻提示"请保持头部稳定"以引导患者自行纠偏；系统同步将整个训练窗口内的头稳采样值连续记录在册。15s窗宽到期后自动终止当前训练，提取窗口内的平均头稳指标作为该轮训练的核心评价参量。

注视训练因其对患者几乎不设主动反应方面的要求——它仅仅要求患者被动地坚守一个固定的头部姿态——因而特别适合作为整个康复训练链条的初始环节，适用于康复早期神经肌肉协同控制能力尚薄弱的患者群体。

\subsection{模式B：扫视训练}

扫视训练模拟了日常生活中视觉注意力在空间目标之间快速跳转的典型场景，目的在于锻炼患者眼球在视野中执行高速、定向且精准的跳跃式再定位动作。训练开始前，系统在标定环节确定的视野有效区间内预生成一张包含8个随机坐标的目标清单，逐次对应墙面上不同的落脚位置。

在每一轮扫视任务中，舵机云台率先将激光束移至当前目标的预设坐标后点亮激光，患者须以最快的速度扫视搜索到光斑所在位置并按下PA2按键予以确认。系统精确记录从激光点亮时刻开始至按键下降沿触发为止的间隔时长，作为本轮扫视的反应耗时。每个目标位置最多驻留3s——如果在此超时窗口内患者未做按键响应，系统将当前轮次标记为超时并播报"时间到咯，再试一次"，随后无条件推进至下一目标。8轮全部执行完毕后，系统汇总统计正确反应次数、平均响应延迟以及正确率三项指标，并据此输出分等级的评价结果。

为提升训练过程中的正向驱动力，本模块额外植入了一条连续正确激励通道：每当患者连续3轮在3s时限内给出有效按键反馈时，语音通道主动播报"非常棒，连续正确，保持状态"作为奖励。

\subsection{模式C：平稳追踪训练}

平稳追踪模式的设计定位是增强患者双眼平滑追随缓慢移动目标的协调能力。早期实现曾采用连续轨迹序列（水平横移→斜向划动→圆形轨道）的方案，然而连续运动模式下缺少在固定航点处的患者确认交互，使得系统无法采集量化反馈以支撑客观评价。当前正式版本已替换为基于离散航点网络（Waypoint-based）的架构。

在标定好的视野空间内预先规划8个均匀分布的航点坐标。每当舵机需要在相邻两个航点之间移动时，调用SmoothStep缓动插值函数驱动云台，使激光光斑以视觉上平滑连续的方式在2s内完成过渡：

\begin{equation}
\text{smoothstep}(t) = t^2 \cdot (3 - 2t),\quad t \in [0, 1]
\label{eq:smoothstep}
\end{equation}

激光到达下一航点后转入检查驻留阶段：光点在当前位置静止保持，患者被要求持续用视线锁住光点并按PA2键确认，系统记录目标到达时刻到按键响应之间的时间差。若患者在2s内未作确认，该航点按超时处理并自动推进至序列中的下一个位置。全部8个航点遍历完成后训练自动终止。

运行期间系统同步监控代偿性转头行为：累计识别到超过10次代偿事件时，在终末评价语音中附加"请尝试减少转头动作"的建议提示；单次代偿事件发生的瞬时则会即时收到"请用眼睛跟踪，不要转头"的提醒。

\subsection{模式D：视觉聚焦训练}

本模式聚焦于睫状肌缩放调节机能的专项锻炼——即双眼在近距离目标与远距离目标之间快速切换焦距的能力。训练场所内的双目标配置为：近距参照点选用头戴装置框架上距患者眼睛约25cm处的一枚红色LED，远距参照点则为投射至前方墙面上的激光光斑。

交替节奏设为5s一个周期：奇数周期内近距LED点亮、远距激光遮断，引导患者将注视焦点收敛在近处；偶数周期内近距LED熄灭、远距激光导通，训练焦点随之拉至墙面。这一近—远交替循环总计往复5轮（即10个阶段），单组训练总时长约50s。

每当患者察觉到墙面激光光斑的出现时按PA2键予以反馈，系统将一次按键视作一个完整的聚焦切换确认事件。该模式不仅可作为独立训练模块使用，也适合在日常完整训练间歇穿插接入，用以调节因长期面对单一刺激模式而引发的视觉疲劳单调感。

\subsection{模式E：空间忽略训练}

空间忽略是脑卒中后较常见的一类认知功能障碍表征，患者的典型特征是对偏盲侧半野内出现的视觉信息加工能力大幅衰减，表现为主体感觉中的"忽略"现象。空间忽略训练的根本目标是通过将视觉刺激与听觉引导信息进行跨通道捆绑，迫使患者主动将注意力资源配送至被长期忽视的半侧空间区域，从而推动双侧空间感知的均衡化重建。

本模式的执行方案采取了左右两侧轮流刺激的设计，每侧各经历3轮试次共计6轮。在每一轮中，舵机先将激光定位至预设的左侧或右侧视野坐标——此坐标由标定阶段测得的X轴边界变量自动运算给出——紧接着语音播报"请注意左/右侧空间"并点亮激光。患者被要求有意识地以视线主动搜索该侧的光点位置，找到后按下PA2按键完成一轮响应，系统记下从点亮激光到按键触发的完整时长。若5s过后患者仍未锁定光点，系统将其判定为超时并播报鼓励语"超时，请多注意该侧空间"。

6轮全部结束后，系统分别计算左侧视野和右侧视野的平均反应延时。患侧与健侧之间的时间差值构成了衡量当前空间忽略严重程度和追踪训练进展的两项客观量化指标。
\section{标定模式}

前述四种训练模式（注视稳定性、扫视、平稳追踪、空间忽略）的实际效果强烈依赖于激光投影位置是否落在每位患者个体化的舒适视野区间之内。患者的头盔佩戴倾角、座椅姿态习惯以及距离墙面的远近各不不同，如若采用固化写入的舵机角度参数，激光点完全有可能漂移至部分患者的有效可视范围之外，导致训练目标不可见或迫使患者以代偿性转头姿态去寻找目标，反而违背了康复初衷。鉴于这一现实需求，系统单独开辟了一条独立的标定通道，允许每次佩戴完成后由患者本人或陪同护理人员借助简单的按键操作对视野的四个边界重新进行个性化测量。

标定流程通过说出"标定模式"命令词激活，系统随即切至舵机标定状态。整个标定过程被拆解为两个串联阶段——X轴（水平）标定先行，Y轴（竖直）标定紧随其后。X轴阶段启动时舵机置Y轴于正中位不动，X轴从较小的起步角开始以恒定速率逐渐增大，激光光斑由此在墙面上自一侧向另一侧平稳扫掠。当患者主观判断光斑已进入个人舒适视野的右边界时按下PA2按键，系统立即抓取并暂存当前角度值calib\_save1；待光斑继续移动至左边界处再次按下按键，将此时的角度存入calib\_save2。X轴双点采集完毕后，X轴回中复位，Y轴标定按相同逻辑接续——Y轴自下端朝上匀速扫描，患者依照前一步的方法依次标记舒适视野的下边界与上边界所对应的舵机输出角度。

标定结果的最终赋值不依赖于按键的先后顺序，而是通过简单的min/max逻辑自动判别边界值。无论患者先按左界还是右界（或先按顶边还是底边），系统始终将较小的角度值存入CALIB\_X\_MIN，较大的角度值存入CALIB\_X\_MAX，Y轴同理。若患者在某个扫描阶段未按键即到达扫描终点，系统自动以默认值填充，避免因漏按键导致的零值异常。标定完成后，四个全局变量CALIB\_X\_MIN、CALIB\_X\_MAX、CALIB\_Y\_MIN和CALIB\_Y\_MAX被更新为本次标定结果，所有训练模式随即自动采用新的视野边界。这些值保存在RAM中，掉电丢失，若需永久保存需重新写入源程序默认值。

\section{语音交互与反馈机制}

\subsection{CI1302语音模块协议}

系统以深圳亚博智能科技基于CI1302芯片构建的AI语音交互模组~\cite{ci1302_doc}作为语音输入输出通道。该模组内嵌了一颗神经网络推理引擎，能够在不依赖外部云端的条件下同时完成离线语音识别（ASR）与文本转语音（TTS）合成两项功能，最大指令容量超过110条，远场拾音半径约为5m。

在硬件连接层面，模组通过STM32的UART7外设（PE7用作接收、PE8用作发送）与主控建立双向数据链路，通信波特率固定为115200bps~\cite{ci1302_doc}。所有物理层数据帧长度统一为5个字节，帧结构定为：$0xAA\ 0x55\ [TYPE]\ [ID]\ 0xFB$，其中$0xAA$和$0x55$充当同步引导头，$0xFB$作为帧尾定界符，中间TYPE字段依据功能分成以下三个域。

表\ref{tab:voice_examples}摘取了各功能域内具有代表性的命令词与播报词。

（1）TYPE=$0x00$：自定义命令词域。当使用者说出某一预设命令词后，模块对其进行本地识别匹配并经由UART回送一条内容相同的5字节应答帧，主控程序以轮询形式调用语音命令获取函数读取返回帧中的ID字节来完成解码。该域内覆盖了全部五种训练模式的启动命令以及标定模式、模式切换等全局控制命令。模式类命令词的设计允许其在空闲等待态、训练执行中、安全暂停态和反馈评价态四个运行阶段被跨状态接受，从而赋予患者在不退出当前训练流程的条件下直接切换训练项目的能力。

（2）TYPE=$0x01\sim0x0A$：唤醒词及相关控制命令域。该域内的指令主要用于训练流程的全局调控，包含暂停、继续、重新开始、跳过当前模式以及唤醒口令"你好小盈"等。此类命令以编码值$0x80 \mid TYPE$的格式经上层统一派发，与控制域TYPE=$0x00$的ID空间相互排他。一旦模块捕捉到"你好小盈"唤醒信号，系统即打开一段持续5s的语音监听窗口，在此时间窗内患者可以直接说出随后的操控指令而无需再次进行唤醒口令的重复呼叫。

（3）TYPE=$0xFF$：纯TTS播报域。主控通过主动向模块发送本域的5字节命令帧来驱动预置语音片段的播放，整个过程中无需等待模块侧的语音识别结果。当前版本共编排了46条独立的TTS播报脚本，ID编码跨度为$0x01$至$0x2D$，按职能划为七个分组——基础反馈组（$0x01\sim0x0B$）、分级评价组（$0x0C\sim0x0E$）、康复激励组（$0x0F\sim0x14$）、过渡衔接组（$0x15\sim0x16$）、指导纠错组（$0x17\sim0x1C$）、专项评价组（$0x1D\sim0x23$）以及训练引导组（$0x24\sim0x2D$）。

\begin{table}[h!]
\centering
\caption{语音交互命令词与播报词示例}
\label{tab:voice_examples}
\begin{tabular}{p{2.5cm} p{3cm} p{6.0cm} p{1.2cm}}
\toprule[1.5pt]
\textbf{协议类型} & \textbf{举例} & \textbf{播报语句} & \textbf{ID} \\
\midrule[0.75pt]
TYPE=$0x00$ & 注视稳定训练 & 注视训练开始，请注视光点保持头部稳定 & 0x02 \\
 命令词区 & 扫视训练 & 扫视训练开始，请快速找到光点并按按键 & 0x03 \\
            & 标定模式 & 进入标定模式，请保持头部居中 & 0x07 \\
            & 模式切换 & 按一下按键进入完整训练模式，按两下进入自选模式 & 0x08 \\
\midrule[0.75pt]
TYPE=$0x01\sim0x0A$ & 暂停训练 & 好的，训练已暂停 & 0x01 \\
 唤醒词区 & 继续训练 & 好的，训练继续 & 0x06 \\
                   & 重新开始 & 好的，重新开始本模式 & 0x09 \\
\midrule[0.75pt]
TYPE=$0xFF$ & 校准完成 & 校准完成 & 0x01 \\
 播报区 & 答得很棒 & 答得很棒 & 0x03 \\
            & 时间到 & 时间到咯，再试一次 & 0x04 \\
            & 训练非常棒 & 表现非常出色，正确率很高，继续加油 & 0x0C \\
            & 全部完成 & 全部训练完成，您辛苦了，休息一下吧 & 0x16 \\
\bottomrule[1.5pt]
\end{tabular}
\end{table}

训练引导组集成了各训练模式专用的开场提示（如注视训练的"注视训练开始，请注视光点保持头部稳定"对应ID $0x27$、扫视训练对应$0x28$等）、标定启动提示、模式选择指引、下一模式确认询问以及完整/自选模式的进入确认语。这批引导播报词为每次训练启程前的语音预报提供了标准化的听觉信号，帮助老年患者在进入体力消耗阶段前建立起明确的行为预期。

模块的UART数据接收完全交由中断事件链驱动：通过调用UART中断接收函数发起逐字节循环接收，每个字节传送完毕即转入串口中断回调钩子，在新字节到达时将回调内集成的帧同步状态机触发一次迁移——该状态机沿序列依次匹配帧同步头$0xAA$和$0x55$，聚齐5字节后校验帧尾$0xFB$，提取中间的TYPE和ID字段并挂起就绪标记。在任意中间阶段检测到字节序列偏离协议格式，状态机便立即复位并转入新一轮帧头搜索状态。
\subsection{命令系统与训练调度}
主控程序在状态机主运行函数的入口处采用了分层命令解析流水线来集中处理各类外部输入。流水线首先轮询式调用语音命令获取函数收取最近一次识别到的命令编码，而后根据编码值的高位域进行路由分发。对于落在唤醒词区间（编码$\ge 0x80$）的指令，系统结合当前所处的运行时状态从中选择匹配的执行分支。其中最特别的是"你好小盈"唤醒触发项：该词被识别后，系统立刻冻结当前正在进行的全部操作动作并将状态强制迁入语音监听暂停态，同步启动5s的倒计时监视器。在此5s的听觉窗口存续期内，患者被允许跳过二次唤醒的冗余步骤而直接发声说出任意后续控制命令——例如"暂停训练""继续训练""重新开始"或是任意一种训练模式的名称——这些在窗口内被捕获到的命令将被立即解码并执行。一旦5s计时耗尽而仍未收到任何跟随指令，系统则自动从暂停中恢复训练。这种压缩式的单唤醒→单命令交互设计消除了传统"唤醒→等待响应→下达命令"三段式对话所引入的时间延滞，对老年患者操控负担的改善尤为明显。



“暂停训练”“继续训练”“重新开始”“跳过这个”等唤醒词区命令在训练状态中具有全局优先级——训练执行中、暂停中或反馈评价中均可通过语音触发，系统按对应的暂停、恢复或模式跳转逻辑执行。

对于TYPE=$0x00$的命令词（模式训练命令和模式切换命令），系统在空闲状态、训练状态、暂停状态和反馈状态四个状态下均予以响应。当用户在训练执行中或暂停状态中说出另一模式名称时，系统调用模式选择与校准函数统一接口——立即停止当前训练、执行陀螺仪校准、然后直接进入新模式的系统训练提示状态流程。这一跨状态模式切换能力允许患者在康复过程中根据自身感受随时灵活调整训练内容，而无需退出当前训练再重新发起。当用户说出“模式切换”命令词时，系统立即进入系统模式选择状态，患者可通过PA2重新选择完整或自选训练流程。

\subsection{分级评价与康复激励}
紧贴老年患者康复进程中的心理接受特点，系统搭建了一套完整的、层级分明的成绩评定与正反馈激励链路。所有TTS播报词均刻意设置了偏慢的朗读语速，每一条播报消息的播放时长被预估并预留出2$\sim$3s的独占窗口，相邻两条播报之间插入了充分的静音间隔以防止语音通道出现消息堆叠与串扰。



在训练结束后，系统根据量化指标进行分级评价：对于扫视训练，根据正确率将评价分为“非常出色”（$\ge80\%$）、“表现不错”（$50\%\sim80\%$）和“还需努力”（$<50\%$）三个等级，同时根据平均反应速度播报专项评价（“反应速度很快”或“反应速度还需提升”）；对于注视训练，根据头稳指标分为“良好”（$<3^\circ$）、“一般”（$3^\circ\sim5^\circ$）和“较差”（$>5^\circ$）三级；对于追踪训练，根据总代偿次数给出“跟踪能力良好”或“尝试减少转头”的专项建议。此外，完成四种训练模式后系统自动汇总并播报“全部训练完成，您辛苦了”。

训练过程中的实时激励包括：扫视训练中连续3次正确响应时的额外表扬提示，以及在训练过半时的进度鼓励播报。

\section{头戴式装置安全保护机制}

安全问题是医疗康复级可穿戴设备设计过程中不可妥协的第一位准则，本系统从三个互为补充的维度夯实了安全保护体系，并将全部保护逻辑与训练状态机的运转深度啮合。

（1）姿态边界监控。主循环在每个调度周期内持续采集俯仰角（roll，对应前后倾）和横滚角（pitch，对应左右转头）的瞬时读数，任一轴的绝对值一旦超出全局安全阈值即置位$safety\_alert$告警标记。与各训练模式内部的自检逻辑有所区分，主循环层的姿态守卫在注视稳定性模式（MODE\_A\_FIXATION）下是刻意被绕过的——此模式本身的任务目标就是让患者自我约束头部晃动，模式内部已通过独立的头稳指标检测（5°阈值超限触发语音提示）形成了专门的行为引导回路，此时若同时在全局层面重复触发暂停机制只会给患者造成不必要的双重告警困扰。

（2）训练过程自动暂停与无损恢复。在训练执行态下，每帧循环均调用姿态安全检查函数对头部姿态做独立校验：一旦捕捉到任一轴的角度值越界，系统立即将状态机迁至安全暂停态。暂停期间，系统完成一组标准化的安全保护动作——关断激光输出与聚焦LED、驱动舵机执行机械回中、触发生蜂鸣器以短促脉冲报警并同步播报语音警示。与此同时，系统持续跟踪姿态的恢复情况，当检测到患者连续3s内全部轴的角度读数均回落至安全区间内时即自动解除暂停并从中断位置无缝接续——舵机角度、训练计数器、倒计时进度等所有上下文变量均原样复原。

（3）全局语音紧急切断。鉴于老年使用者可能在训练期间的任意时刻出现不适或需要求助的特殊情况，系统部署了一条覆盖全部状态的语音紧急中断链路。无论当前处于训练执行、安全暂停还是反馈评价等任意阶段，患者只需说出唤醒口令"你好小盈"，CI1302模组完成本地识别后立刻回传对应编码，主控在收到信号的下一帧内立即执行全系统静默——激光熄灭、舵机归零位——并将状态跳转至空闲等待状态挂起等候。该操作的独到之处在于：执行过程中不清除任何训练相关的进度变量，患者若在短时间内确认状况恢复，仅需说出"继续训练"即可从打断位置直接接续，完全免去了重新经历模式选择和设备校准等前置步骤的冗余时间损耗。
